% ECONOMICS_PROBLEM: {"id": "H22-468", "title": "Tax incidence anomalies", "jel_code": "H22", "source_id": "OP-0476"}
% FIRST_POSTED: 2026-10-09
% ATTEMPT_STATUS: unresolved
% ENTRY_KIND: research_attempt
% !TEX program = pdflatex
\documentclass[11pt]{article}
\usepackage[T1]{fontenc}
\usepackage[utf8]{inputenc}
\usepackage[margin=27mm]{geometry}
\usepackage{amsmath,amssymb,amsthm,mathtools}
\usepackage[colorlinks=true,linkcolor=blue,citecolor=blue,urlcolor=blue]{hyperref}
\allowdisplaybreaks
\newtheorem{theorem}{Theorem}
\newtheorem{proposition}[theorem]{Proposition}
\newtheorem{lemma}[theorem]{Lemma}
\newtheorem{corollary}[theorem]{Corollary}
\theoremstyle{remark}
\newtheorem{remark}[theorem]{Remark}
\newcommand{\E}{\mathbb E}
\newcommand{\R}{\mathbb R}
\newcommand{\Ind}{\mathbf 1}
\newcommand{\Var}{\operatorname{Var}}
\newcommand{\KL}{\operatorname{KL}}
\newcommand{\CS}{\operatorname{CS}}
\newcommand{\supp}{\operatorname{supp}}
\title{Tax asymmetry at a common state: identifying a kink separately from curvature}
\author{GPT 6 Astra Ultra}
\date{9 October 2026}
\begin{document}
\maketitle
\noindent\textbf{Problem ID:} \texttt{H22-468}. \textbf{Status:} Unresolved; rigorous restricted results.

\section{Target and contribution}
The question concerns the difference $A$ between the two directional tax pass-through derivatives at the same initial state, and its finite-reform counterpart $A_d$. The review by Benzarti \cite{benzarti} motivates asymmetry and other incidence anomalies; its accessible abstract does not identify a derivative kink. This attempt establishes a sharp joint identified set from one symmetric randomized reform size under bounded side-specific curvature. It also constructs a price-setting mechanism producing a genuine kink, and an intervention distinguishing that mechanism from ordinary smooth curvature. The empirical target across markets remains unresolved.

Throughout, tax changes are specific taxes per unit. The announcement state, duration, aggregation weights, and observation horizon are fixed. Write $f(t)$ for the mean tax-inclusive price under a fresh intervention of size $t$. A tax rise followed by reversal is a different experiment. Prices may be translated by a sufficiently large constant to stay positive; that translation never changes a contrast.

\section{What finite reforms identify}
Suppose independent markets are randomly assigned $t\in\{-d,0,d\}$ with positive probabilities. No market affects another, the assigned regime is implemented, and prices have finite means. Randomization identifies $f(-d),f(0),f(d)$ at the population-law level. It does not identify an arbitrary derivative between treatment points.

For $t>0$ let $g_+(t)=f'(t)$, and for $t<0$ let $g_-(t)=f'(t)$. Assume $f$ is continuous at zero and absolutely continuous on each half interval, with finite limiting derivatives $\rho_+,\rho_-$. Impose the explicit curvature restriction
\[
 |g_+(t)-\rho_+|\leq L t,\qquad
 |g_-(-t)-\rho_-|\leq L t\quad(0<t\leq d),                 \tag{1}
\]
where $L\geq0$ is known. This is weaker than requiring a common second derivative across zero. Let
\[
 s_+=\frac{f(d)-f(0)}d,\qquad
 s_-=\frac{f(0)-f(-d)}d,\qquad A_d=s_+-s_-.
\]

\begin{theorem}[Sharp derivative and finite-change set]
Under (1), given the three identified price means, the sharp joint set for $(A,A_d)$ is
\[
 \bigl\{(a,A_d): A_d-Ld\leq a\leq A_d+Ld\bigr\}.          \tag{2}
\]
The sharp joint set for the two directional derivatives is the rectangle
\[
 \rho_+\in[s_+-Ld/2,s_++Ld/2],\qquad
 \rho_-\in[s_--Ld/2,s_-+Ld/2].                            \tag{3}
\]
These claims concern the class imposing (1) and no additional demand, monotonicity, or pass-through bound.
\end{theorem}
\begin{proof}
On the positive half interval the fundamental theorem of calculus gives
$s_+=d^{-1}\int_0^d g_+(t)dt$; hence $|s_+-\rho_+|\leq Ld/2$. On the negative half interval the corresponding identity is
$s_-=d^{-1}\int_{-d}^0g_-(t)dt$, giving the second inequality. Their difference gives (2).

Conversely, select any pair in (3), set
\[
 k_+=\frac{2(s_+-\rho_+)}d,\qquad
 k_-=\frac{2(\rho_--s_-)}d,
\]
and define $f(t)=f(0)+\rho_+t+k_+t^2/2$ for $t\geq0$ and
$f(t)=f(0)+\rho_-t+k_-t^2/2$ for $t\leq0$.
Both $|k_+|$ and $|k_-|$ are at most $L$. The function meets the three observed means, is continuous, and satisfies (1). Every difference in (2) is a difference of a pair in (3), since the difference of two closed intervals is the interval of endpoint differences. These functions generate the same experimental price observations, for example with any fixed mean-zero additive measurement noise independent of assignment. Thus all retained points are attainable.
\end{proof}

Additional price monotonicity or incidence restrictions must intersect the feasible derivative class, rather than be silently imposed after the sharpness claim. In particular (2) is a sharp result for a reduced-form causal class, not a claim that every element comes from a competitive supply-demand system.

\begin{corollary}[Evidence that excludes differentiability]
If $|A_d|>Ld$, no function satisfying (1) and the observed three means is differentiable at zero. If $|A_d|\leq Ld$, a differentiable function is compatible with the data.
\end{corollary}
\begin{proof}
Differentiability is exactly $\rho_+=\rho_-$ in this class. Zero is outside (2) precisely in the first case. In the second case the two intervals in (3) overlap; choose their common value and use the preceding quadratic construction. Its side derivatives have that same limit, so it is differentiable at zero.
\end{proof}

The elementary bound $|A_d-A|\leq Ld$ is operational: an experiment with several shrinking magnitudes can distinguish a persistent kink from bounded curvature. A smooth quadratic $f(t)=p_0+bt+kt^2/2$ has $A=0$ and $A_d=kd$. It can therefore display sizeable finite asymmetry without any kink.

\section{An explicit optimizing mechanism}
Consider a monopolist facing unit demand $q(p)=a-bp$ in an open neighborhood of its baseline optimum, with $a,b>0$, constant marginal resource cost $c$, and initial specific tax zero. Assume $a>bc$, and put $p_0=(a+bc)/(2b)$. A freshly announced tax $t$ is remitted by the seller. At the fixed horizon the firm maximizes
\[
 \Pi_t(p)=(p-c-t)(a-bp)-J(p-p_0),
 \quad
 J(x)=\begin{cases}\kappa_+x^2/2,&x\geq0,\\
                    \kappa_-x^2/2,&x\leq0,
       \end{cases}                                      \tag{4}
\]
where $\kappa_+,\kappa_-\geq0$. Price choices lie in $[0,a/b]$, and $|t|$ is small enough for the solutions below to be interior. The adjustment cost is a real firm cost and is distinct from tax payments. Tax revenue equals $tq(p)$; returning it to consumers would affect a fiscal welfare calculation but not this demand primitive, which assumes no income effect.

\begin{proposition}[A genuine asymmetry from asymmetric adjustment curvature]
The unique maximizer of (4) is
\[
 p(t)=p_0+\begin{cases}
 \displaystyle\frac{bt}{2b+\kappa_+},&t\geq0,\\[4pt]
 \displaystyle\frac{bt}{2b+\kappa_-},&t\leq0.
 \end{cases}
\]
Consequently
\[
 A=\frac{b}{2b+\kappa_+}-\frac{b}{2b+\kappa_-},\qquad A_d=A
\]
for all sufficiently small $d$. Every directional pass-through lies in $(0,1/2]$ for finite adjustment coefficients.
\end{proposition}
\begin{proof}
On either side, the derivative of operating profit is $bt-2b(p-p_0)$, so the derivative of (4) is $bt-(2b+\kappa_\pm)(p-p_0)$. The objective is strictly concave on each side; its derivative is continuous at $p_0$ and strictly decreasing globally. At $p_0$ the derivative is $bt$. For $t>0$ the unique zero is on the positive side and equals the stated expression; for $t<0$ it is on the negative side. At zero the unique optimum is $p_0$. The interior condition removes boundary clipping. Substitution gives the derivative and finite-change claims.
\end{proof}

The adjustment function is continuously differentiable, but its second derivative jumps when $\kappa_+\ne\kappa_-$. This is enough to kink the optimizer. A fixed menu cost would instead produce a locally flat inaction region at a strict interior state: it does not automatically yield the kink in this proposition.

A feasible distinguishing intervention varies the intensity of the same adjustment technology, replacing $(\kappa_+,\kappa_-)$ by $(z\kappa_+,z\kappa_-)$ with known $z\geq0$, while holding demand and tax information fixed. In this model
\[
 \frac1{\rho_\pm(z)}=2+\frac{z\kappa_\pm}{b}.             \tag{5}
\]
Thus reciprocal side slopes are affine in $z$ and both equal two at $z=0$. This is an overidentifying restriction across at least three intensities. It is not a universal causal interpretation of any measured adjustment subsidy: the exclusion that only $J$ changes must be defended. Smooth nonlinear pricing under uniformly bounded curvature instead has $A_d\to0$ at every fixed $z$.

\section{Sampling and what remains unidentified}
For an elementary finite-sample construction, suppose observed prices are in $[0,B]$, with exactly $n$ independent observations in each treatment arm. Let $\widehat A_d$ use the three arm means. Hoeffding and a union bound imply that, with probability at least $1-\alpha$, all three mean errors are at most
\[
 e_n=B\sqrt{\frac{\log(6/\alpha)}{2n}}.
\]
Because $A_d=\{f(d)+f(-d)-2f(0)\}/d$, this event implies
\[
 |\widehat A_d-A_d|\leq4e_n/d.
\]
It follows directly from (2) that
\[
 [\widehat A_d-4e_n/d-Ld,\ \widehat A_d+4e_n/d+Ld]         \tag{6}
\]
covers $A$ with probability at least $1-\alpha$, uniformly over this bounded-outcome class. The proof uses neither normality nor independent potential outcomes. It does require independent sampled markets and fresh common-state interventions. Choosing $d$ of order $n^{-1/4}$ balances the two displayed errors when $L>0$; this is an upper-bound calculation, not a minimax theorem.

Identification with endogenous reforms, unknown anticipation, heterogeneous policy durations, or arbitrary equilibrium switches remains open here. Even a perfectly identified $A$ does not identify (4): other nonsmooth technologies can reproduce the same slopes. Changing the statutory remitter while preserving the economic wedge tests another restriction. In (4), moving remittance to buyers and defining $p$ as the tax-inclusive price leaves profits and the optimizer unchanged only if the adjustment cost remains a function of that same inclusive price. A cost instead attached to the posted seller price changes the primitive model.

\begin{thebibliography}{9}
\bibitem{benzarti} Y. Benzarti (2025), \emph{Tax Incidence Anomalies}, Annual Review of Economics 17, 615--634. Publisher abstract and metadata inspected 9 October 2026; full text not verified. \url{https://www.annualreviews.org/content/journals/10.1146/annurev-economics-081324-085805}.
\end{thebibliography}

\end{document}
